\section*{अध्याय \ref{ch:g9:alternator} --- बिजलीघर से घर तक: प्रत्यावर्तित्र}

\begin{solution}{exo:g9:alternator:1}
प्रेरण सिर्फ़ बदलाव करता है; ज़्यादा तेज़ हरकत, ज़्यादा तगड़ा \omterm{def:g1:magnets:magnet}{चुंबक} और
ज़्यादा फेरे बड़ा \omterm{def:g8:voltage-voltmeter:voltage}{विभवांतर} देते हैं; और उलटी हरकत धारा उलट देती है।
ठहरा हुआ \omterm{def:g1:magnets:magnet}{चुंबक}, चाहे कितना ही तगड़ा हो, ठीक-ठीक कुछ भी प्रेरित नहीं
करता।
\end{solution}

\begin{solution}{exo:g9:alternator:2}
घूर्णक का हर आधा चक्कर एक ध्रुव को कुंडली की ओर लाता है (एक ओर प्रेरण),
और फिर उसे दूर ले जाता है (दूसरी ओर): इसलिए \omterm{def:g8:voltage-voltmeter:voltage}{विभवांतर} को हर आधे घूर्णन
पर अपना चिह्न बदलना ही पड़ता है --- प्रत्यावर्तन उस घुमाव में ही बना
हुआ है।
\end{solution}

\begin{solution}{exo:g9:alternator:3}
हर सेकंड पचास चक्कर (सादी दो-ध्रुव वाली मशीन के लिए)। और हर बिजलीघर को
ठीक वही घूर्णन-लय थामे रखनी पड़ती है --- पूरा जाल एक क़दम पर घूमता है
और एक ही \omterm{def:g8:sound-pitch-loudness:frequency}{आवृत्ति} बाँटता है।
\end{solution}

\begin{solution}{exo:g9:alternator:4}
टरबाइन और \omterm{def:g9:alternator:alternator}{प्रत्यावर्तित्र}। बदलता सिर्फ़ पंखुड़ियों को धकेलने वाला है:
लपट की भाप, परमाणु की भाप, गिरती झील, या \omterm{def:g2:air-around-us:air}{हवा}।
\end{solution}

\begin{solution}{exo:g9:alternator:5}
झील की जमा ऊँचाई $\to$ गिरते पानी की गति $\to$ टरबाइन का घूर्णन $\to$
\omterm{def:g9:alternator:alternator}{प्रत्यावर्तित्र} की बिजली $\to$ ट्रांसफ़ॉर्मर और जाल $\to$ लैंप। और वह
झील सूर्य ने भरी: वाष्पीकरण ने समुद्र उठाया, और बारिश ने उसे पहाड़ों
तक पहुँचा दिया।
\end{solution}

\begin{solution}{exo:g9:alternator:6}
लैंप की \omterm{def:g5:everyday-energy:energy}{ऊर्जा} कहीं से तो आनी ही है, और साइकिल पर अकेला भंडार सवार ही
है: डाइनेमो टाँगों की मेहनत को \omterm{def:g1:light-and-shadows:source}{प्रकाश} में बदल देता है, इसलिए लैंप के
\omterm{def:g9:kinetic-energy-safety:joule}{जूल} फ़ालतू खिंचाव बनकर सामने आते हैं। \omterm{def:g5:everyday-energy:energy}{ऊर्जा} बदली जाती है, जादू से पैदा
कभी नहीं की जाती --- पाठ्यक्रम का सबसे पुराना नियम।
\end{solution}

\begin{solution}{exo:g9:alternator:7}
\omterm{ex:g9:alternator:chains}{नवीकरणीय}: \omterm{def:g2:air-around-us:air}{हवा}, झील (बारिश उसे दोबारा भर देती है), धूप। नहीं: कोयला,
यूरेनियम --- जो इंसानी पैमानों पर हमेशा के लिए ख़र्च हो जाते हैं। छँटाई
यही पूछती है: जितना हम निकालते हैं, क्या क़ुदरत उतना भर देती है?
\end{solution}

\begin{solution}{exo:g9:alternator:8}
$P = U I$ एक ही \omterm{def:g9:electric-power-energy:power}{शक्ति} को ऊँचे \omterm{def:g8:voltage-voltmeter:voltage}{विभवांतर} पर छोटी धारा बनकर सफ़र करने देता
है; और \omterm{prop:g8:electric-current-effects:three}{ऊष्मीय प्रभाव} तारों का नुक़सान धारा के साथ बढ़ा देता है। छोटी
धारा, ठंडे तार, बची हुई \omterm{def:g5:everyday-energy:energy}{ऊर्जा} --- इसलिए जाल ऊँचे पर ढोता है और देहरी
पर उतार देता है।
\end{solution}

\begin{solution}{exo:g9:alternator:9}
धीमा घुमाव: छोटा प्रेरित \omterm{def:g8:voltage-voltmeter:voltage}{विभवांतर} (दूसरा नियम) --- मद्धम लैंप --- और
हर सेकंड चंद चक्र: यानी आँख के मिलाने की दर से नीचे, साफ़ टिमटिमाहट।
तेज़ घुमाव: बड़ा \omterm{def:g8:voltage-voltmeter:voltage}{विभवांतर} --- चमकीला --- और ऐसी \omterm{def:g8:sound-pitch-loudness:frequency}{आवृत्ति} जिसे आँख
चिकनाकर मिला देती है। दो घुंडियाँ, एक ही हत्था।
\end{solution}

\begin{solution}{exo:g9:alternator:10}
\qty{230}{V} पर: $I = \num{1e8} \div 230 \approx \qty{4.3e5}{A}$ ---
यानी चार लाख \omterm{def:g8:intensity-ammeter:intensity}{ऐम्पियर}: पृथ्वी पर ऐसी कोई केबल नहीं। और \qty{400000}{V}
पर: $I = \num{1e8} \div \num{4e5} = \qty{250}{A}$ --- यानी आम संचरण
केबल। ट्रांसफ़ॉर्मर का मुक़दमा बंद।
\end{solution}

\begin{solution}{exo:g9:alternator:11}
दोनों शृंखलाएँ यूँ पढ़ी जाती हैं: भंडार $\to$ \omterm{def:g5:heat-and-insulation:heat}{ऊष्मा} $\to$ उबलता पानी
$\to$ भाप की गति $\to$ घूर्णन $\to$ बिजली --- टरबाइनों को खिलाती
विस्तृत केतलियाँ। और वे ठीक एक ही कड़ी पर अलग होती हैं: \omterm{def:g5:heat-and-insulation:heat}{ऊष्मा} बनाता
कौन है --- जलते कोयले का रसायन, या टूटते यूरेनियम के नाभिक।
\end{solution}

\begin{solution}{exo:g9:alternator:12}
उलटी दिशा में: दीवार की \qty{230}{V} मोहल्ले के ट्रांसफ़ॉर्मर से आई
(\omterm{def:g8:voltage-voltmeter:voltage}{विभवांतर} उतारा हुआ), वह मीनारों पर लाखों \omterm{def:g8:voltage-voltmeter:voltage}{वोल्ट} पर चली (ढोने का रूप:
ऊँचे \omterm{def:g8:voltage-voltmeter:voltage}{विभवांतर} वाली AC, जो थोड़ी गरमाहट रिसाती है), और वह किसी
\omterm{def:g9:alternator:alternator}{प्रत्यावर्तित्र} से आई (प्रेरण से घूर्णन को बिजली में बदलते हुए), जिसे
किसी टरबाइन ने घुमाया (भाप या पानी की गति से), और जिसे किसी भंडार ने
खिलाया --- मान लो बाँध की उठाई हुई झील ने, जिसे सूर्य के वाष्पीकरण ने
उठाया था। तो नौवीं \omterm{def:g6:sun-earth-moon:axisorbit}{कक्षा} वाला वाक्य: तुम्हारे कमरे का \omterm{def:g1:light-and-shadows:source}{प्रकाश} धूप ही है
--- जमा की हुई, गिराई हुई, घुमाई हुई और ढोकर लाई हुई।
\end{solution}

\begin{solution}{pb:g9:alternator:1}
\textbf{1.} झील की ऊँचाई $\to$ पानी की झपटती गति $\to$ टरबाइन का
घूर्णन $\to$ \omterm{def:g9:alternator:alternator}{प्रत्यावर्तित्र} का प्रेरित प्रत्यावर्ती \omterm{def:g8:voltage-voltmeter:voltage}{विभवांतर} $\to$
बिजलीघर का बसबार।
\textbf{2.} घूर्णक वाला \omterm{def:g1:magnets:magnet}{चुंबक} घूमता है; और कुंडलियाँ ठहरी रहती हैं।
हरकत के बिना कुंडलियों में से गुज़रते चुंबकत्व में कोई बदलाव ही नहीं
होगा --- और बिना बदलाव के प्रेरण कुछ भी नहीं देता।
\textbf{3.} \omterm{def:g8:voltage-voltmeter:voltage}{विभवांतर} की \omterm{def:g8:sound-pitch-loudness:frequency}{आवृत्ति} घूर्णन की नक़ल करती है: ग़लत \omterm{def:g7:motion-average-speed:formula}{चाल},
ग़लत \omterm{def:g8:sound-pitch-loudness:frequency}{आवृत्ति} --- और जाल की \qty{50}{Hz} से क़दम-बेक़दम हुई मशीन जाल के
बाक़ी हर \omterm{def:g9:alternator:alternator}{प्रत्यावर्तित्र} की मदद करने के बजाय उनसे लड़ने लगती है।
\textbf{4.} घूर्णक का ज़्यादा तगड़ा \omterm{def:g1:magnets:magnet}{चुंबक} सचमुच ज़्यादा प्रेरित करता
है --- और उसी पल टरबाइन को उस कड़े चुंबकीय खिंचाव के ख़िलाफ़ उसे
घुमाते रहने के लिए ज़्यादा मेहनत करनी पड़ती है: ज़्यादा बिजली बाहर
ज़्यादा पानी भीतर माँगती है, \omterm{def:g9:kinetic-energy-safety:joule}{जूल} दर \omterm{def:g9:kinetic-energy-safety:joule}{जूल}। दाम झील चुकाती है, \omterm{def:g1:magnets:magnet}{चुंबक}
नहीं।
\textbf{5.} $I = \num{2e7} \div \num{2e5} = \qty{100}{A}$.
\textbf{6.} $I = \num{2e7} \div 230 \approx \qty{8.7e4}{A}$ ---
यानी सत्तासी हज़ार \omterm{def:g8:intensity-ammeter:intensity}{ऐम्पियर}: नामुमकिन केबल। ट्रांसफ़ॉर्मर का काम ठीक
यही अदला-बदली है --- \omterm{def:g8:voltage-voltmeter:voltage}{विभवांतर} ऊपर, धारा नीचे, और \omterm{def:g9:electric-power-energy:power}{शक्ति} जस की तस।
\textbf{7.} फाटकों से ज़्यादा पानी। कम खिलाने पर बोझ तले दबी मशीनें
चढ़ाई पर चढ़ते साइकिल सवार की तरह धीमी पड़ जाती हैं --- और जाल की
\omterm{def:g8:sound-pitch-loudness:frequency}{आवृत्ति} पचास से नीचे ढल जाती है: जाल की भूख उसकी हर घड़ी और हर मशीन
महसूस करती है।
\textbf{8.} सिर्फ़ पहला डिब्बा --- यानी \omterm{def:g2:air-around-us:air}{हवा}, पंखुड़ियों का धक्का।
टरबाइन, \omterm{def:g9:alternator:alternator}{प्रत्यावर्तित्र} और ट्रांसफ़ॉर्मर पूरी तरह भले-चंगे खड़े रहे,
बस बेदम।
\textbf{9.} जाल एक विशाल समांतर जाल है: हर बिजलीघर को हर क़स्बे से कई
सड़कें जोड़ती हैं, इसलिए कटी हुई लाइन का हिस्सा बची हुई \omterm{def:g5:series-parallel-circuits:parallel}{शाखाओं} से बह
सकता है --- समांतर की स्वतंत्रता, राष्ट्रीय पैमाने पर।
\textbf{10.} घुमाई गई, अतिभारित लाइनों ने एक मिनट के लिए पहुँचाया गया
\omterm{def:g8:voltage-voltmeter:voltage}{विभवांतर} नामिक से नीचे ढलने दिया: गोदाम के लैंप मंद पड़ गए और उसकी
मोटरें आलस से घूमीं --- यानी भूखी, जैसा अपने \omterm{def:g8:voltage-voltmeter:nominal}{नामिक विभवांतर} से नीचे
रखे हर उपकरण को होना ही है।
\textbf{11.} ``आज रात का हर \omterm{def:g9:electric-power-energy:power}{वाट} गिरती झील के पानी से आया; और समुद्र को
सूर्य ने वाष्प बनाया तथा बारिश ने उसे हमारे पहाड़ों तक पहुँचाया। हम
यहाँ धूप ही घुमाते हैं --- जो आकाश के रास्ते एक झील में जमा हुई थी।''
\textbf{12.} जैसे: ``चुंबकों के पास से गुज़रती गति ने रात का हर \omterm{def:g8:voltage-voltmeter:voltage}{वोल्ट}
बनाया --- और सात बजे दस हज़ार केतलियों ने झील से ज़रा और तेज़ गिरने
को कह दिया।''
\end{solution}
